\subsection{Remark on modular weights in low-energy effective field theory from type II string theory --- arXiv:2301.10356}

\textbf{Authors:} Shota Kikuchi, Tatsuo Kobayashi, Kaito Nasu, Hajime Otsuka, Shohei Takada, Hikaru Uchida

\paragraph{主な主張}
磁化D-brane模型の物質場のモジュラー重みは局在磁束や曲率によって補正され得る一方、物理的Yukawa結合は変わらず、超対称性破れのsfermion質量には差が生じ得ると示す~\cite{Kikuchi:2023clx}。

\paragraph{新規性}
単純な零モード波動関数から得るモジュラー重みを局在効果まで含めて再検討し、物理的結合とsoft質量への影響を区別する。

\paragraph{位置付け}
弦理論由来のモジュラー重みを低エネルギー模型で採用する際に、その幾何学的補正と物理的帰結を検討する研究。
